% Compilar desde este directorio con:
% pdflatex Clase_05_Diseno_experimental_de_Machine_Learning.tex
\documentclass[aspectratio=169,10pt]{beamer}

\usepackage[utf8]{inputenc}
\usepackage[T1]{fontenc}
\usepackage[spanish,es-nodecimaldot]{babel}
\usepackage{amsmath,amssymb,booktabs,tabularx,array,tikz}
\usetikzlibrary{arrows.meta,positioning,calc,fit,shapes.geometric}

\definecolor{UAPNavy}{HTML}{17324D}
\definecolor{UAPTeal}{HTML}{007F82}
\definecolor{UAPGold}{HTML}{E2A33A}
\definecolor{UAPLight}{HTML}{F3F6F8}
\definecolor{UAPGray}{HTML}{536471}
\definecolor{UAPRed}{HTML}{A94343}

\tikzset{
  flow/.style={-{Stealth[length=2.2mm]},thick,draw=UAPNavy},
  datum/.style={draw=UAPGold,fill=UAPGold!16,rounded corners=2pt,align=center,minimum height=8mm},
  process/.style={draw=UAPTeal,fill=UAPTeal!12,rounded corners=3pt,very thick,align=center,minimum height=10mm},
  decision/.style={draw=UAPNavy,fill=UAPNavy!9,rounded corners=3pt,very thick,align=center,minimum height=10mm}
}

\usetheme{Boadilla}
\usecolortheme{default}
\setbeamercolor{normal text}{fg=UAPNavy,bg=white}
\setbeamercolor{structure}{fg=UAPTeal}
\setbeamercolor{frametitle}{fg=white,bg=UAPNavy}
\setbeamercolor{title}{fg=white}
\setbeamercolor{subtitle}{fg=UAPGold}
\setbeamercolor{block title}{fg=white,bg=UAPTeal}
\setbeamercolor{block body}{fg=UAPNavy,bg=UAPLight}
\setbeamercolor{alertblock title}{fg=white,bg=UAPRed}
\setbeamercolor{alertblock body}{fg=UAPNavy,bg=UAPRed!7}
\setbeamerfont{frametitle}{series=\bfseries}
\setbeamertemplate{navigation symbols}{}
\setbeamertemplate{itemize item}{\color{UAPGold}$\blacktriangleright$}
\setbeamertemplate{itemize subitem}{\color{UAPTeal}$\bullet$}
\setbeamertemplate{footline}{%
  \leavevmode\hbox{%
    \begin{beamercolorbox}[wd=.78\paperwidth,ht=2.8ex,dp=1.2ex,leftskip=1em]{author in head/foot}%
      \color{white}\insertshorttitle
    \end{beamercolorbox}%
    \begin{beamercolorbox}[wd=.22\paperwidth,ht=2.8ex,dp=1.2ex,rightskip=1em plus1fil]{date in head/foot}%
      \color{UAPGray}\hfill Semana 5 \enspace | \enspace \insertframenumber/\inserttotalframenumber
    \end{beamercolorbox}%
  }%
}

\newcommand{\concepto}[1]{\textcolor{UAPTeal}{\textbf{#1}}}
\newcommand{\explicacion}[1]{\par\vspace{.25em}{\scriptsize\color{UAPGray}\textit{#1}\par}}

\title[Data Science · Clase 5]{Diseño experimental de Machine Learning}
\subtitle{Cómo formular, evaluar y comparar modelos de manera confiable}
\author{Curso de Data Science}
\institute{Semana 5 · Clase 5}
\date{}

\begin{document}

% 1
\begin{frame}[plain]
  \begin{beamercolorbox}[wd=\paperwidth,ht=\paperheight,sep=2em]{frametitle}
    \vfill
    {\usebeamerfont{title}\usebeamercolor[fg]{title}\Huge\inserttitle\par}
    \vspace{.45em}
    {\usebeamerfont{subtitle}\usebeamercolor[fg]{subtitle}\Large\insertsubtitle\par}
    \vspace{1.1em}
    \begin{tikzpicture}[node distance=4.5mm,scale=.72,transform shape,every node/.style={text=UAPNavy}]
      \node[datum] (p) {Problema}; \node[process,right=of p] (t) {Tarea};
      \node[datum,right=of t] (s) {Partición}; \node[process,right=of s] (e) {Evaluación};
      \node[decision,right=of e] (r) {Evidencia};
      \draw[flow] (p)--(t); \draw[flow] (t)--(s); \draw[flow] (s)--(e); \draw[flow] (e)--(r);
    \end{tikzpicture}\par
    \vspace{.7em}
    {\color{white}\small Un resultado es confiable cuando el experimento representa el uso futuro y evita información indebida.\par}
    \vfill
    {\color{UAPGold}\large\insertauthor\par}
    {\color{white}\insertinstitute\par}
  \end{beamercolorbox}
\end{frame}

% 2
\begin{frame}{De la incertidumbre al aprendizaje}
  \small
  En la Semana 4 estudiamos estimaciones, variabilidad y decisiones. Ahora usaremos datos históricos para construir una función que produzca estimaciones en casos nuevos.
  \vspace{.4em}
  \begin{tabularx}{\textwidth}{X X}
    \toprule Antes & Ahora \\ \midrule
    Estimar una proporción o una media & Aprender una relación entre variables \\
    Cuantificar variabilidad muestral & Medir desempeño fuera del ajuste \\
    Comparar acciones y costos & Elegir métricas vinculadas con esos costos \\
    \bottomrule
  \end{tabularx}
  \begin{block}{Continuidad}
    Una métrica también es una estimación obtenida sobre una muestra y bajo supuestos determinados.
  \end{block}
\end{frame}

% 3
\begin{frame}{Propósito y resultados de aprendizaje}
  \footnotesize
  \textbf{Propósito:} diseñar una evaluación de Machine Learning que produzca evidencia relevante, reproducible y libre de fuga de información.
  \begin{block}{Al finalizar, podremos}
    \begin{itemize}
      \item distinguir paradigma, tarea, modelo, predicción y decisión;
      \item formular clasificación, regresión, clustering y series temporales;
      \item elegir particiones coherentes con la población de uso;
      \item establecer baselines y métricas pertinentes;
      \item aplicar validación cruzada sin contaminar la evaluación;
      \item reconocer subajuste, sobreajuste y comparaciones injustas;
      \item documentar un protocolo experimental reproducible.
    \end{itemize}
  \end{block}
\end{frame}

% 4
\begin{frame}{Pregunta de apertura}
  \small
  Una organización desea anticipar problemas de calidad del agua para priorizar inspecciones.
  \begin{itemize}
    \item ¿Qué debe predecirse exactamente?
    \item ¿Qué información estará disponible al predecir?
    \item ¿Se evaluarán nuevas mediciones, nuevos sitios o meses futuros?
    \item ¿Qué error tiene mayor costo?
    \item ¿Cómo sabremos que un modelo mejora la práctica actual?
  \end{itemize}
  \begin{block}{Idea}
    Antes de elegir un algoritmo debemos definir qué evidencia demostraría que su resultado es útil fuera de los datos usados para construirlo.
  \end{block}
\end{frame}

% 5
\begin{frame}{¿Qué significa aprender de los datos?}
  \small
  Aprender es ajustar una función a partir de ejemplos:
  \[
    f:\mathcal X\rightarrow\mathcal Y,\qquad \widehat y=f(x).
  \]
  \begin{itemize}
    \item $x$: variables disponibles para un caso.
    \item $y$: resultado real que se desea estimar.
    \item $\widehat y$: predicción producida por el modelo.
    \item El ajuste utiliza casos históricos.
    \item La evaluación debe usar casos que no participaron del ajuste.
  \end{itemize}
  \begin{alertblock}{Resultado}Memorizar el pasado no demuestra que el modelo generalice.\end{alertblock}
\end{frame}

% 6
\begin{frame}{Paradigma, tarea y algoritmo}
  \small
  \begin{tabularx}{\textwidth}{>{\bfseries}p{.18\textwidth} X X}
    \toprule Nivel & Pregunta & Ejemplo \\ \midrule
    Paradigma & ¿De dónde proviene la señal de aprendizaje? & supervisado o no supervisado \\
    Tarea & ¿Qué salida necesitamos? & clase, valor, grupo o pronóstico \\
    Algoritmo & ¿Qué procedimiento ajustará la relación? & regresión, árbol o K-means \\
    \bottomrule
  \end{tabularx}
  \vspace{.6em}
  Una misma tarea admite varios algoritmos y un algoritmo puede resolver tareas diferentes.
  \explicacion{Elegir ``un árbol'' no define todavía la población, la etiqueta, la partición ni la métrica.}
\end{frame}

% 7
\begin{frame}{Aprendizaje supervisado}
  \small
  El conjunto de entrenamiento contiene pares con respuesta conocida:
  \[
    \mathcal D=\{(x_1,y_1),\ldots,(x_n,y_n)\}.
  \]
  \begin{itemize}
    \item La respuesta $y$ guía el ajuste.
    \item En \concepto{clasificación}, $y$ representa una categoría.
    \item En \concepto{regresión}, $y$ representa una cantidad.
    \item El desempeño se mide comparando $\widehat y$ con respuestas reservadas.
  \end{itemize}
  \begin{block}{Ejemplo}
    Variables fisicoquímicas y antecedentes permiten estimar si una futura medición excederá un límite.
  \end{block}
\end{frame}

% 8
\begin{frame}{Aprendizaje no supervisado y semisupervisado}
  \small
  \textbf{No supervisado:} trabaja con entradas $x$ sin una respuesta objetivo conocida.
  \begin{itemize}
    \item Busca estructura, similitudes o representaciones.
    \item El clustering agrupa sitios con perfiles semejantes.
    \item Los grupos requieren interpretación y validación.
  \end{itemize}
  \textbf{Semisupervisado:} combina pocos casos etiquetados con muchos casos sin etiqueta.
  \begin{itemize}
    \item Puede ser útil cuando medir es barato y confirmar el resultado es costoso.
    \item La evaluación final sigue necesitando etiquetas confiables.
  \end{itemize}
\end{frame}

% 9
\begin{frame}{Cuatro tareas de análisis}
  \footnotesize
  \begin{tabularx}{\textwidth}{>{\bfseries}p{.18\textwidth} p{.27\textwidth} X}
    \toprule Tarea & Salida & Pregunta de evaluación \\ \midrule
    Clasificación & categoría o probabilidad & ¿detecta correctamente las clases? \\
    Regresión & valor numérico & ¿cuánto se aleja del valor real? \\
    Clustering & grupo o pertenencia & ¿los grupos son estables e interpretables? \\
    Serie temporal & valor futuro & ¿supera un pronóstico ingenuo respetando el tiempo? \\
    \bottomrule
  \end{tabularx}
  \begin{alertblock}{Control}
    La codificación numérica de una categoría no la convierte en regresión.
  \end{alertblock}
\end{frame}

% 10
\begin{frame}{Un dominio, cuatro formulaciones}
  \small
  Con datos de calidad del agua podemos formular problemas distintos:
  \begin{tabularx}{\textwidth}{X p{.2\textwidth} X}
    \toprule Necesidad & Tarea & Objetivo \\ \midrule
    Activar una inspección & clasificación & excede/no excede un límite \\
    Estimar concentración & regresión & concentración en mg/L \\
    Comparar perfiles de sitios & clustering & agrupación por similitud \\
    Anticipar el próximo mes & serie temporal & concentración futura \\
    \bottomrule
  \end{tabularx}
  \begin{block}{Resultado}
    Cada formulación cambia la unidad, la salida, la partición y el criterio de evaluación.
  \end{block}
\end{frame}

% 11
\begin{frame}{Definición del problema predictivo}
  \small
  Antes de modelar debemos establecer:
  \begin{enumerate}
    \item decisión o uso de la predicción;
    \item unidad de análisis y población objetivo;
    \item instante en que se predice y horizonte del resultado;
    \item variables disponibles;
    \item objetivo y fuente de confirmación;
    \item criterio de evaluación.
  \end{enumerate}
  \begin{block}{Pregunta de control}
    ¿Qué cambiará operativamente cuando el modelo produzca una predicción?
  \end{block}
\end{frame}

% 12
\begin{frame}{Unidad, población y muestra}
  \small
  \begin{itemize}
    \item \concepto{Unidad:} entidad sobre la que se produce una predicción; por ejemplo, \texttt{(sitio, fecha)}.
    \item \concepto{Población objetivo:} conjunto de unidades donde se pretende utilizar el modelo.
    \item \concepto{Muestra histórica:} unidades observadas disponibles para ajustar y evaluar.
  \end{itemize}
  \begin{alertblock}{Dependencia}
    La misma estación en dos fechas produce dos filas, pero esas filas pueden estar relacionadas.
  \end{alertblock}
  \begin{block}{Consecuencia}
    El número de filas no siempre coincide con el número de unidades independientes.
  \end{block}
\end{frame}

% 13
\begin{frame}{Instante y horizonte}
  \small
  \begin{itemize}
    \item $t_0$: instante en que debe emitirse la predicción.
    \item Ventana de observación: información disponible hasta $t_0$.
    \item Horizonte $H$: periodo futuro donde se define el resultado.
    \item Confirmación: momento en que la etiqueta puede conocerse.
  \end{itemize}
  \vspace{.4em}\centering
  \begin{tikzpicture}[node distance=6mm,scale=.75,transform shape]
    \node[datum] (p) {Pasado observable}; \node[process,right=of p] (t) {$t_0$: predicción};
    \node[decision,right=of t] (h) {Horizonte $H$}; \node[datum,right=of h] (r) {Resultado};
    \draw[flow] (p)--(t); \draw[flow] (t)--(h); \draw[flow] (h)--(r);
  \end{tikzpicture}
  \explicacion{Cambiar $t_0$ o $H$ cambia el problema, aunque las columnas tengan los mismos nombres.}
\end{frame}

% 14
\begin{frame}{Generalización}
  \small
  Generalizar es mantener un desempeño útil en casos nuevos de la población objetivo.
  \[
    R(f)=\mathbb E_{(X,Y)\sim P_{\mathrm{uso}}}[L(Y,f(X))].
  \]
  La evaluación aproxima ese error esperado con datos no utilizados para ajustar el procedimiento.
  \begin{itemize}
    \item Casos representativos del uso esperado.
    \item Separación compatible con tiempo y entidades.
    \item Métrica alineada con el error relevante.
  \end{itemize}
  \begin{block}{Alcance}El resultado vale para el escenario evaluado, no para cualquier población futura.\end{block}
\end{frame}

% 15
\begin{frame}{Entrenamiento, validación y prueba}
  \small
  \begin{tabularx}{\textwidth}{>{\bfseries}p{.22\textwidth} X}
    \toprule Partición & Uso permitido \\ \midrule
    Entrenamiento & ajustar parámetros y transformaciones \\
    Validación & elegir variables, configuración y umbral \\
    Prueba & medir una vez el procedimiento congelado \\
    \bottomrule
  \end{tabularx}
  \begin{alertblock}{Control}
    Consultar repetidamente la prueba y modificar el procedimiento convierte la prueba en validación.
  \end{alertblock}
  \explicacion{Los porcentajes 60/20/20 u 80/10/10 son opciones, no reglas universales.}
\end{frame}

% 16
\begin{frame}{Partición aleatoria y estratificada}
  \small
  \begin{itemize}
    \item \concepto{Aleatoria:} asigna unidades al azar.
    \item \concepto{Estratificada:} conserva aproximadamente la proporción de clases.
    \item Son adecuadas cuando las unidades pueden tratarse como intercambiables e independientes.
    \item La estratificación estabiliza la medición de clases poco frecuentes.
  \end{itemize}
  \begin{alertblock}{Limitación}
    Ninguna evita que mediciones del mismo sitio o del mismo episodio aparezcan a ambos lados.
  \end{alertblock}
\end{frame}

% 17
\begin{frame}{Partición por grupos}
  \small
  Todos los registros de un grupo permanecen en una sola partición.
  \begin{block}{Ejemplo}
    Sitios A--F $\rightarrow$ entrenamiento \quad|\quad G--H $\rightarrow$ validación \quad|\quad I--J $\rightarrow$ prueba.
  \end{block}
  \begin{itemize}
    \item Grupo posible: sitio, sensor, paciente, cliente, vehículo o evento.
    \item Evita compartir patrones propios de una entidad.
    \item Evalúa transferencia hacia grupos no observados.
  \end{itemize}
  \textbf{Pregunta:} ¿el uso real predice nuevas fechas de sitios conocidos o sitios completamente nuevos?
\end{frame}

% 18
\begin{frame}{Partición temporal}
  \small
  El pasado entrena y el futuro evalúa:
  \begin{block}{Orden}
    enero--junio: entrenamiento $\rightarrow$ julio--agosto: validación $\rightarrow$ septiembre: prueba.
  \end{block}
  \begin{itemize}
    \item Respeta el orden causal y simula un despliegue prospectivo.
    \item Revela cambios de estación, sensores o procesos.
    \item Cada transformación se ajusta solo con datos anteriores al periodo evaluado.
  \end{itemize}
  \begin{alertblock}{Riesgo}
    Barajar observaciones temporales puede permitir que el futuro ayude a predecir el pasado.
  \end{alertblock}
\end{frame}

% 19
\begin{frame}{¿Qué partición representa el uso?}
  \footnotesize
  \begin{tabularx}{\textwidth}{X X}
    \toprule Uso futuro & Partición principal \\ \midrule
    nuevas filas independientes & aleatoria o estratificada \\
    nuevas observaciones de grupos conocidos & temporal dentro de cada grupo \\
    grupos nunca observados & por grupos \\
    meses futuros & temporal \\
    nuevos grupos en periodos futuros & combinación de grupos y tiempo \\
    \bottomrule
  \end{tabularx}
  \begin{block}{Criterio}
    Primero se define qué significa ``caso nuevo''; después se elige la partición.
  \end{block}
\end{frame}

% 20
\begin{frame}{Fuga de información}
  \small
  Existe \concepto{fuga} cuando el ajuste o la evaluación utilizan información que no estaría disponible al producir una predicción real.
  \begin{itemize}
    \item Produce resultados demasiado optimistas.
    \item Puede aparecer en variables, particiones o transformaciones.
    \item También ocurre cuando unidades relacionadas cruzan la frontera experimental.
  \end{itemize}
  \begin{block}{Prueba operativa}
    Para cada variable registrar quién la genera, cuándo existe y qué datos intervinieron en su cálculo.
  \end{block}
\end{frame}

% 21
\begin{frame}{Tres fugas frecuentes}
  \footnotesize
  \begin{tabularx}{\textwidth}{>{\bfseries}p{.18\textwidth} X X}
    \toprule Tipo & Ejemplo & Control \\ \midrule
    Temporal & usar la medición que confirma el evento & cortar variables en $t_0$ \\
    Entidad & mismo sitio en entrenamiento y prueba & separar por grupo \\
    Preprocesamiento & imputar con la media de todos los datos & ajustar solo con entrenamiento \\
    \bottomrule
  \end{tabularx}
  \begin{alertblock}{Resultado}
    Eliminar una columna sospechosa no basta: toda la frontera experimental debe preservar la separación.
  \end{alertblock}
\end{frame}

% 22
\begin{frame}{Ejemplo numérico: imputación contaminada}
  \small
  Valores de entrenamiento: $[8,10,12]$. Valor faltante en entrenamiento: \texttt{NA}. Valor de prueba: $30$.
  \begin{itemize}
    \item Media correcta, calculada solo con entrenamiento:
    \[
      \bar x_{train}=(8+10+12)/3=10.
    \]
    \item Media contaminada, calculada incluyendo prueba:
    \[
      \bar x_{todo}=(8+10+12+30)/4=15.
    \]
  \end{itemize}
  \begin{alertblock}{Problema}
    Imputar con 15 permite que el valor reservado de prueba modifique los datos de entrenamiento.
  \end{alertblock}
\end{frame}

% 23
\begin{frame}{Duplicados y selección de atributos}
  \small
  \textbf{Duplicados o casos relacionados:} copias, ventanas solapadas o mediciones del mismo episodio pueden inflar la muestra y aparecer en particiones distintas.
  \vspace{.5em}

  \textbf{Selección de atributos:} elegir variables mediante correlaciones, rankings o desempeño también aprende de los datos.
  \begin{itemize}
    \item Los duplicados relacionados deben permanecer juntos.
    \item La selección basada en datos debe realizarse dentro de entrenamiento o de cada fold.
    \item La prueba no interviene en ninguna selección.
  \end{itemize}
\end{frame}

% 24
\begin{frame}{Baseline: referencia mínima}
  \small
  Un baseline es una referencia simple, reproducible y disponible con la misma información que el modelo.
  \begin{itemize}
    \item Determina qué mejora debe justificar la complejidad.
    \item Detecta errores del pipeline y métricas mal interpretadas.
    \item Puede ser una regla operativa existente, no solo una constante.
    \item Debe evaluarse en exactamente los mismos casos.
  \end{itemize}
  \begin{block}{Resultado}
    Superar un baseline débil no demuestra utilidad; la referencia debe ser pertinente.
  \end{block}
\end{frame}

% 25
\begin{frame}{Baselines por tipo de tarea}
  \small
  \begin{tabularx}{\textwidth}{>{\bfseries}p{.22\textwidth} X}
    \toprule Tarea & Baseline posible \\ \midrule
    Clasificación & clase mayoritaria o frecuencia histórica \\
    Regresión & media o mediana del entrenamiento \\
    Serie temporal & último valor o mismo periodo anterior \\
    Clustering & partición trivial y estabilidad frente a remuestreo \\
    \bottomrule
  \end{tabularx}
  \begin{block}{Referencia operativa}
    En calidad del agua también puede compararse contra la regla de inspección vigente.
  \end{block}
\end{frame}

% 26
\begin{frame}{Ejemplo: accuracy engañosa}
  \footnotesize
  En 100 mediciones hay 20 excedencias y 80 casos normales. El baseline siempre predice ``normal''.
  \begin{center}
  \begin{tabular}{lrr}
    \toprule & Excedencia real & Normal real \\ \midrule
    Predice excedencia & VP = 0 & FP = 0 \\
    Predice normal & FN = 20 & VN = 80 \\
    \bottomrule
  \end{tabular}
  \end{center}
  \[
    \mathrm{accuracy}=80/100=80\%,\qquad \mathrm{recall}=0/20=0\%.
  \]
  \begin{alertblock}{Lectura}
    Tiene accuracy alta, pero no detecta ninguna excedencia.
  \end{alertblock}
\end{frame}

% 27
\begin{frame}{Matriz de confusión}
  \small
  \begin{center}
  \begin{tabular}{lcc}
    \toprule & Real positiva & Real negativa \\ \midrule
    Predicción positiva & verdadero positivo (VP) & falso positivo (FP) \\
    Predicción negativa & falso negativo (FN) & verdadero negativo (VN) \\
    \bottomrule
  \end{tabular}
  \end{center}
  \begin{itemize}
    \item VP y VN son aciertos respecto de la clase definida como positiva.
    \item FP es una falsa alarma; FN es un positivo omitido.
    \item La clase positiva debe declararse antes de interpretar las métricas.
  \end{itemize}
\end{frame}

% 28
\begin{frame}{Métricas de clasificación}
  \small
  Supongamos $VP=16$, $FP=8$, $FN=4$ y $VN=72$:
  \[
    \mathrm{accuracy}=\frac{16+72}{100}=0{,}88,
  \]
  \[
    \operatorname{Prec}=\frac{16}{16+8}=0{,}67,\qquad
    \mathrm{recall}=\frac{16}{16+4}=0{,}80,
  \]
  \[
    F_1=\frac{2VP}{2VP+FP+FN}=\frac{32}{44}=0{,}73.
  \]
  \begin{block}{Resultado}Cada métrica responde una pregunta diferente.\end{block}
\end{frame}

% 29
\begin{frame}{Métrica, prevalencia y costo}
  \small
  \begin{itemize}
    \item \concepto{Accuracy:} fracción total de aciertos; puede ocultar una clase rara.
    \item \concepto{Precisión:} de las alertas emitidas, cuántas eran correctas.
    \item \concepto{Recall:} de los eventos reales, cuántos se detectaron.
    \item \concepto{$F_1$:} equilibra precisión y recall, pero no incorpora verdaderos negativos.
  \end{itemize}
  Si omitir una excedencia es más costoso que realizar una inspección innecesaria, recall merece especial atención.
  \begin{alertblock}{Control}
    Ninguna métrica sustituye la declaración de costos y capacidad operativa.
  \end{alertblock}
\end{frame}

% 30
\begin{frame}{Métricas de regresión}
  \small
  Para errores $e_i=y_i-\widehat y_i$:
  \[
    \mathrm{MAE}=\frac{1}{n}\sum_{i=1}^n|e_i|,
    \qquad
    \mathrm{RMSE}=\sqrt{\frac{1}{n}\sum_{i=1}^n e_i^2}.
  \]
  \begin{itemize}
    \item Ambas conservan la unidad de $y$.
    \item MAE pondera linealmente la magnitud del error.
    \item RMSE penaliza más los errores grandes.
    \item La elección depende de las consecuencias de los errores extremos.
  \end{itemize}
\end{frame}

% 31
\begin{frame}{Ejemplo numérico: MAE y RMSE}
  \small
  Valores reales: $[10,12,18,20]$. Predicciones: $[11,10,17,26]$.

  Errores absolutos: $[1,2,1,6]$.
  \[
    \mathrm{MAE}=\frac{1+2+1+6}{4}=2{,}5.
  \]
  \[
    \mathrm{RMSE}=\sqrt{\frac{1^2+2^2+1^2+6^2}{4}}
    =\sqrt{10{,}5}\approx3{,}24.
  \]
  \begin{block}{Lectura}
    RMSE aumenta más por el error de 6 unidades.
  \end{block}
\end{frame}

% 32
\begin{frame}{$R^2$ y sus límites}
  \small
  \[
    R^2=1-\frac{\sum_i(y_i-\widehat y_i)^2}{\sum_i(y_i-\bar y)^2}.
  \]
  \begin{itemize}
    \item Compara el error del modelo con predecir la media.
    \item Puede ser negativo fuera de muestra.
    \item No conserva la unidad del objetivo.
    \item Un valor alto no garantiza errores operativamente pequeños.
    \item En series con tendencia puede resultar engañoso.
  \end{itemize}
  \begin{block}{Recomendación}
    Informar $R^2$ junto con MAE o RMSE y un baseline pertinente.
  \end{block}
\end{frame}

% 33
\begin{frame}{Evaluar clustering}
  \small
  Sin etiquetas objetivo no existe una accuracy predictiva convencional.
  \begin{itemize}
    \item \concepto{Cohesión:} observaciones próximas dentro de cada grupo.
    \item \concepto{Separación:} grupos alejados entre sí.
    \item \concepto{Silhouette:} combina cohesión y separación en una escala aproximada de $-1$ a $1$.
    \item \concepto{Estabilidad:} grupos similares ante remuestreo o cambios razonables.
    \item \concepto{Utilidad:} perfiles interpretables que apoyan una decisión.
  \end{itemize}
  \begin{alertblock}{Control}
    Un valor interno alto no demuestra que los grupos sean relevantes para el dominio.
  \end{alertblock}
\end{frame}

% 34
\begin{frame}{Evaluar series temporales}
  \small
  \begin{itemize}
    \item Entrenar con pasado y validar sobre periodos posteriores.
    \item Comparar contra persistencia y estacionalidad simple.
    \item Medir varios orígenes y horizontes, no un único corte conveniente.
    \item Ajustar transformaciones nuevamente en cada ventana.
  \end{itemize}
  \begin{block}{Ventana expansiva}
    enero--marzo $\rightarrow$ abril \quad|\quad enero--abril $\rightarrow$ mayo \quad|\quad enero--mayo $\rightarrow$ junio.
  \end{block}
  \explicacion{El promedio resume desempeño; la variación entre periodos revela estabilidad temporal.}
\end{frame}

% 35
\begin{frame}{Validación cruzada}
  \small
  En K-fold, el conjunto de desarrollo se divide en $K$ bloques. Cada bloque actúa una vez como validación:
  \[
    \bar m=\frac{1}{K}\sum_{k=1}^{K}m_k.
  \]
  \begin{itemize}
    \item Estratificada para conservar clases.
    \item Por grupos para mantener entidades juntas.
    \item Temporal con ventanas ordenadas.
    \item La prueba final permanece fuera del ciclo.
  \end{itemize}
  \begin{block}{Resultado}
    Reportar promedio y dispersión entre folds evita depender de una única partición favorable.
  \end{block}
\end{frame}

% 36
\begin{frame}{El pipeline se ajusta dentro de cada fold}
  \footnotesize
  En cada iteración:
  \begin{enumerate}
    \item ajustar imputación, escalado y codificación con los folds de entrenamiento;
    \item seleccionar atributos usando solo esos folds;
    \item transformar entrenamiento y validación;
    \item ajustar el modelo;
    \item medir sobre el fold reservado.
  \end{enumerate}
  \vspace{.25em}\centering
  \begin{tikzpicture}[node distance=4mm,scale=.66,transform shape]
    \node[datum] (t) {Training del fold}; \node[process,right=of t] (p) {Transformaciones};
    \node[process,right=of p] (s) {Selección}; \node[decision,right=of s] (m) {Modelo};
    \node[datum,right=of m] (v) {Validación};
    \draw[flow] (t)--(p); \draw[flow] (p)--(s); \draw[flow] (s)--(m); \draw[flow] (m)--(v);
  \end{tikzpicture}
  \explicacion{El objeto evaluado es el procedimiento completo, no solo el algoritmo final.}
\end{frame}

% 37
\begin{frame}{Subajuste y sobreajuste}
  \footnotesize
  \begin{tabularx}{\textwidth}{>{\bfseries}p{.18\textwidth} p{.2\textwidth} p{.2\textwidth} X}
    \toprule Situación & Entrenamiento & Validación & Lectura \\ \midrule
    Subajuste & error alto & error alto & capacidad o representación insuficiente \\
    Equilibrio & error bajo & error bajo y estable & generalización plausible \\
    Sobreajuste & error muy bajo & error mayor & adaptación excesiva al entrenamiento \\
    \bottomrule
  \end{tabularx}
  \begin{block}{Criterio}
    Se elige por desempeño de validación, estabilidad y simplicidad, no por el mejor ajuste al entrenamiento.
  \end{block}
\end{frame}

% 38
\begin{frame}{Actividad: protocolo experimental}
  \scriptsize
  Cada equipo definirá, antes de entrenar:
  \begin{tabularx}{\textwidth}{>{\bfseries}p{.18\textwidth} X}
    \toprule Sección & Evidencia mínima \\ \midrule
    Problema & unidad, población, $t_0$, horizonte, $x$ e $y$ \\
    Partición & entrenamiento, validación y prueba justificadas \\
    Controles & riesgos de fuga y cómo se comprobarán \\
    Referencia & baseline reproducible \\
    Evaluación & métrica principal, auxiliares y costo relevante \\
    Reproducción & versión de datos, semilla y reglas congeladas \\
    \bottomrule
  \end{tabularx}
  \begin{block}{Hito}
    Otro equipo debe poder reconstruir y auditar la evaluación propuesta.
  \end{block}
\end{frame}

% 39
\begin{frame}{Síntesis y próximo paso}
  \footnotesize
  \textbf{Flujo:} definición $\rightarrow$ tarea $\rightarrow$ partición $\rightarrow$ baseline $\rightarrow$ pipeline $\rightarrow$ validación $\rightarrow$ prueba.
  \begin{itemize}
    \item La tarea determina la salida, no el algoritmo.
    \item La partición debe representar qué significa un caso nuevo.
    \item La fuga invalida la evidencia de generalización.
    \item Cada tarea necesita baselines y métricas pertinentes.
    \item La validación evalúa todo el procedimiento.
    \item La prueba mide una vez una configuración congelada.
  \end{itemize}
  \begin{block}{Próximo paso}
    Implementar transformaciones y atributos ajustados únicamente con entrenamiento.
  \end{block}
\end{frame}

\end{document}
